% ==============================================================================
% Fichier : front/00_page_titre.tex
% Titre   : Page de titre - \titreCours
% Auteur  : MINKA MI NGUIDJOÏ Thierry Emmanuel
% Version : 1.1 -- novembre 2026 (charte des cours, modèle de HE2)
% ==============================================================================
% Logos : Université de Garoua à gauche et à droite.
% ==============================================================================

\begin{titlepage}
\thispagestyle{empty}

\begin{tikzpicture}[remember picture, overlay]
    \fill[CouleurPrimaire]
        (current page.north west) rectangle
        ([yshift=-3.6cm] current page.north east);
    \fill[CouleurAccent]
        ([yshift=-3.6cm] current page.north west) rectangle
        ([yshift=-4.0cm] current page.north east);
    \fill[CouleurPrimaire]
        (current page.south west) rectangle
        ([yshift=2.5cm] current page.south east);
    \fill[CouleurAccent]
        ([yshift=2.5cm] current page.south west) rectangle
        ([yshift=2.9cm] current page.south east);
    % En-tête institutionnel, posé dans le bandeau supérieur
    \node[align=center, text=white] at ([yshift=-1.8cm] current page.north) {%
        {\fontsize{11}{13}\selectfont
        \textsc{République du Cameroun}}\\[0.3ex]
        {\color{CouleurAccent!60!white}\rule{0.4\paperwidth}{0.4pt}}\\[0.6ex]
        {\fontsize{10}{12}\selectfont\textbf{\tutelle}}\\[0.2ex]
        {\fontsize{9}{11}\selectfont \faculte{} · \departement}\\[0.2ex]
        {\fontsize{8}{10}\selectfont \unitepedago}};
    % Logos : l'Université de Garoua à gauche et à droite (un seul établissement)
    \node[fill=white, rounded corners=4pt, inner sep=2pt]
        at ([xshift=1.75cm, yshift=-1.8cm] current page.north west)
        {\includegraphics[height=2.6cm]{figures/logos/garoua_logo.png}};
    \node[fill=white, rounded corners=4pt, inner sep=2pt]
        at ([xshift=-1.75cm, yshift=-1.8cm] current page.north east)
        {\includegraphics[height=2.6cm]{figures/logos/garoua_logo.png}};
    % Année et version, posées dans le bandeau inférieur
    \node[align=center] at ([yshift=1.25cm] current page.south) {%
        {\color{white}\bfseries\large Année académique \anneescolaire}\\[0.3ex]
        {\color{CouleurAccent!70!white}\small
        Version 1.0 | Document pédagogique de l'UE \codeUE}};
\end{tikzpicture}

\vspace*{0.6cm}

% --- Titre principal (sur deux lignes) ---
\begin{center}
    {\color{CouleurPrimaire}\rule{0.85\textwidth}{1.2pt}}\\[1.2ex]
    {\fontsize{24}{30}\selectfont\bfseries\color{CouleurPrimaire}
        Systèmes de Gestion de la\\[0.4ex]
        Sécurité Informatique}\\[0.8ex]
    {\color{CouleurAccent}\rule{0.85\textwidth}{2pt}}\\[0.5ex]
    {\color{CouleurPrimaire}\rule{0.85\textwidth}{0.5pt}}
\end{center}

\vspace{0.8cm}

\begin{center}
    {\fontsize{13}{16}\selectfont\itshape\color{CouleurSecondaire}
        Du risque à la certification~:\\
        piloter la sécurité de l'information}
\end{center}

\vspace{0.5cm}

% --- Cartouche sigle et code UE ---
\begin{center}
    \begin{tikzpicture}
        \node[draw=CouleurPrimaire, fill=CouleurPrimaire, rounded corners=6pt,
              inner xsep=14pt, inner ysep=6pt]{%
            \color{white}\bfseries\large \sigleCours{} - UE~\codeUE
        };
    \end{tikzpicture}
\end{center}

\vspace{1.0cm}

% --- Schéma de la discipline : le cycle d'amélioration continue du SMSI ---
\begin{center}
\begin{tikzpicture}[
    f/.style = {draw=CouleurBordDef, fill=CouleurDef, rounded corners=3pt,
        minimum width=2.2cm, minimum height=0.8cm, align=center,
        font=\small\bfseries\color{CouleurPrimaire}},
    arrow/.style = {-{Stealth[length=4pt]}, thick, color=CouleurSecondaire}]
\foreach \t [count=\i] in {Planifier, Déployer, Contrôler, Améliorer}
    \node[f] (n\i) at (3.3*\i - 3.3, 0) {\t};
\foreach \i [evaluate=\i as \j using int(\i+1)] in {1,2,3}
    \draw[arrow] (n\i) -- (n\j);
\draw[arrow] (n4.south) to[bend left=18] node[below, font=\footnotesize\itshape,
    text=CouleurSecondaire] {amélioration continue, guidée par les risques} (n1.south);
\end{tikzpicture}
\end{center}

\vspace{0.6cm}

\begin{center}
    {\large\bfseries\color{CouleurPrimaire}
        \auteurcours}\\[0.5ex]
    {\normalsize\color{CouleurSecondaire}
       Enseignant Chercheur | \tutelle}\\[0.2ex]
    {\small\color{CouleurBordInfo}
        \faculte{} · \departement}\\[0.2ex]
    {\small\color{CouleurBordInfo}
        \href{mailto:maletyon@proton.me}{maletyon@proton.me}}
\end{center}

\end{titlepage}
